\documentclass[letterpaper]{article} % DO NOT CHANGE THIS
\usepackage[submission]{aaai2027}  % DO NOT CHANGE THIS
\usepackage[hyphens]{url}  % DO NOT CHANGE THIS
\usepackage{graphicx} % DO NOT CHANGE THIS
\urlstyle{rm} % DO NOT CHANGE THIS
\def\UrlFont{\rm}  % DO NOT CHANGE THIS
\usepackage{natbib}  % DO NOT CHANGE THIS AND DO NOT ADD ANY OPTIONS TO IT
\usepackage{caption} % DO NOT CHANGE THIS AND DO NOT ADD ANY OPTIONS TO IT
\frenchspacing  % DO NOT CHANGE THIS
\usepackage{amsmath}
\usepackage{amssymb}
\usepackage{amsthm}
\usepackage{algorithm}
\usepackage{algorithmic}
\usepackage{booktabs}
\usepackage{xspace}

\newtheorem{proposition}{Proposition}
\newtheorem{theorem}{Theorem}
\newtheorem{lemma}{Lemma}
\newtheorem{corollary}{Corollary}
\newtheorem{definition}{Definition}
\newtheorem{remark}{Remark}

\pdfinfo{
/TemplateVersion (2027.1)
}

\setcounter{secnumdepth}{2}

\title{Supplementary Material:\\EigenTree: a Centralizing Communication Protocol\\for Decentralized Multi-Agent Multi-Armed Bandit}
\author{
    Anonymous Submission
}
\affiliations{
    %
}

\begin{document}

\nocopyright

\maketitle

Notation follows the main paper throughout. $G=(V,E)$ is the underlying
graph, $N=|V|$, and $A$ its adjacency matrix. $\tilde\psi$ is the
(unnormalized, locally computed) leading eigenvector of $A$, with eigenvalue
$\lambda$, and $h^\star=\arg\max_i\tilde\psi_i$ is the hub. $\mathcal{T}$ is
the induced routing tree and $D$ its height. $\mathrm{par}(i)$, $\mathrm{ch}(i)$, and
$d_i$ denote a node's parent, children, and depth in $\mathcal{T}$, and
$\mathcal{N}_G(i)$ denotes node $i$'s neighbor set in $G$.

\section{Hub Emergence, Routing, and Coordinator Sufficiency}

\begin{proposition}[Centralized aggregation and broadcast via routing]
\label{prop:aggbcast}
Suppose each node $i$ holds a scalar $x_i$ and propagates it toward
$h^\star$ one hop per round, with each intermediate node forwarding the sum
it received. Then the hub holds $\sum_i x_i$ after $D$ rounds, where $D$ is
the tree height. Symmetrically, a value broadcast from $h^\star$ one hop
per round reaches every node after $D$ rounds. Neither direction requires
any node to know $D$ in advance: the uplink terminates when the hub
observes a round with no incoming packets, and each node infers its own
depth $d_i$ from the hop count carried on the first downlink it receives.
\end{proposition}

\begin{proof}
\textbf{Uplink.} We show by induction on depth, from the leaves upward,
that every node at depth $d$ has forwarded the sum of its entire subtree's
values to its parent by the end of round $D-d+1$.

\emph{Base case.} A leaf $\ell$ at depth $d_\ell$ has no children, so its
subtree sum is simply $x_\ell$. It has nothing to wait for, so it forwards
$x_\ell$ to $\mathrm{par}(\ell)$ in round 1.

\emph{Inductive step.} Suppose every node at depth $d+1$ has forwarded its
subtree sum to its parent by the end of round $D-(d+1)+1 = D-d$. A node $i$
at depth $d$ with children $\mathrm{ch}(i)$ waits until it has received a
forwarded value from every child. Each child's value arrives no later than
round $D-d$ by the inductive hypothesis. Node $i$ sums these values
together with its own $x_i$ and forwards the total to $\mathrm{par}(i)$ in
round $D-d+1$.

Since $h^\star$ is at depth 0, the induction applied down to $d=1$ shows
every direct child of $h^\star$ has forwarded its subtree sum by round $D$.
The hub receives all of these by round $D$ and sums them with its own
$x_{h^\star}$ to obtain $\sum_i x_i$.

\emph{Termination without knowing $D$.} A node $i$ at depth $d$ cannot know
$D$, but it does not need to: it simply waits until every child, a fixed
and locally known set, has sent its value, then forwards. The hub likewise
does not know $D$ in advance. It accumulates incoming packets each round
and declares the uplink complete once a round passes with no new arrivals.
Every node forwards exactly once, by the induction above, so the hub
receives no further packets after round $D$, and this silence is a purely
local observation.

\textbf{Downlink.} A symmetric induction on depth, from the root downward,
gives the same result. $h^\star$ sends $v$ to $\mathrm{ch}(h^\star)$ in
round 1, tagged with hop count $0$. A node at depth $d$ receiving $(v,
\mathrm{hops}=d-1)$ sets its own depth $d_i \leftarrow (d-1)+1 = d$, which
is correct since it is exactly $d$ hops from $h^\star$ along tree edges,
and forwards $(v, \mathrm{hops}=d)$ to its own children in the next round.
By induction, every node at depth $d$ receives $v$ in round $d$, so every
node receives $v$ within $D$ rounds. Each node's depth is derived purely
from the hop count it personally receives, with no global knowledge
required.
\end{proof}

\begin{theorem}[Coordinator sufficiency]
\label{thm:coordinator}
Let $\mathcal{A}$ be any decision rule whose output at each step depends
only on globally aggregated per-agent statistics $B = \sum_{i \in V} b_i$,
where $b_i$ is agent $i$'s local contribution. Suppose a coordinator
receives the exact aggregate $B$ within $D$ rounds and broadcasts its
decision within $D$ further rounds, with all agents synchronized via the
depth signal carried on the broadcast. Then the coordinator's information
state at every decision point is identical to that of a centralized
instance of $\mathcal{A}$, with an additive overhead of at most $2D$ time
steps per decision round.
\end{theorem}

\begin{proof}
Fix any decision round $r$. By Proposition~\ref{prop:aggbcast} (uplink
direction, with $x_i = b_i$), after $D$ rounds the hub holds $B =
\sum_{i\in V} b_i$ exactly, not an approximation or a partial sum, since
every node's contribution is included by the termination argument above.
Because $\mathcal{A}$'s decision at round $r$ depends only on $B$, the hub
computes exactly the action a centralized instance of $\mathcal{A}$ would
compute given the same set of per-agent contributions.

By Proposition~\ref{prop:aggbcast} (broadcast direction), within a further
$D$ rounds every agent holds the hub's decision, and by the same mechanism
every agent has derived its own depth $d_i$ from the hop count on this
broadcast. The depth signal is carried on every downlink, so all agents can
identify the round at which the current cycle ends, namely $2D$ rounds
after the cycle's pull phase, and begin the next round's contribution
simultaneously, regardless of their individual depth in $\mathcal{T}$.

Across the whole run, the sequence of aggregates $B^{(1)}, B^{(2)},
\ldots$ seen by the hub is therefore identical, round for round, to the
sequence a centralized instance of $\mathcal{A}$ would see given the same
sequence of per-agent local contributions, so the sequence of decisions is
identical as well. The only difference from a literal centralized run is
the $2D$ rounds of communication latency inserted between successive
decisions: an additive overhead per round that does not alter the sequence
of decisions itself.
\end{proof}

\section{Cost of Not Using the Minimum Diameter Tree}
\label{sec:mdst-gap}

Throughout this section, $D = \mathrm{ecc}_G(h^\star)$ denotes the height
of the EigenTree routing tree, equivalently the eccentricity of $h^\star$ in
$G$. $D' = R(G) = \min_{v\in V}\mathrm{ecc}_G(v)$ denotes the radius of
$G$: the height achieved by rooting a shortest-path tree at $G$'s absolute
center, a minimum diameter spanning tree (MDST) \citep{hakimi1964,hassintamir1995}.

\begin{proposition}[Hub degree lower bound]
\label{prop:hubdeg}
The hub's degree is at least the graph's spectral radius: $\deg(h^\star) \geq \lambda$.
\end{proposition}

\begin{proof}
$\tilde\psi$ is the (unnormalized) leading eigenvector of $A$ with
eigenvalue $\lambda$, so it satisfies $A\tilde\psi = \lambda\tilde\psi$.
Reading off the $h^\star$-th row of this equation,
\[
\lambda\tilde\psi_{h^\star} = \sum_{j \sim h^\star} \tilde\psi_j,
\]
where the sum ranges over the $\deg(h^\star)$ neighbors of $h^\star$ in
$G$. $\tilde\psi_{h^\star} = \max_k \tilde\psi_k$ by definition of
$h^\star$, so every term in the sum satisfies $\tilde\psi_j \leq
\tilde\psi_{h^\star}$, giving
\[
\lambda\tilde\psi_{h^\star} = \sum_{j\sim h^\star}\tilde\psi_j \leq \deg(h^\star)\cdot\tilde\psi_{h^\star}.
\]
By the Perron--Frobenius theorem, the leading eigenvector of a connected
graph's adjacency matrix is strictly positive, so $\tilde\psi_{h^\star} >
0$. We divide both sides by $\tilde\psi_{h^\star}$ to obtain $\lambda \leq
\deg(h^\star)$.
\end{proof}

\begin{theorem}[Dense graphs: depth at most 2]
\label{thm:mdst-gap}
Let $\delta_{\bar N}(h^\star) = \min_{u \notin \mathcal{N}_G(h^\star)\cup\{h^\star\}} \deg(u)$ be the smallest degree among the nodes $h^\star$ is not directly connected to. If $\deg(h^\star) + \delta_{\bar N}(h^\star) \geq N$, then $D \leq 2$.
\end{theorem}

\begin{proof}
Let $u \in V$ be any node with $u \neq h^\star$ and $u \notin
\mathcal{N}_G(h^\star)$: $u$ is not $h^\star$ itself and not a neighbor of
$h^\star$. We show $d(h^\star, u) \leq 2$. Since $u$ was an arbitrary
non-neighbor, and every neighbor of $h^\star$ trivially has
$d(h^\star,\cdot)=1$, this establishes $D = \mathrm{ecc}_G(h^\star) \leq 2$.

Consider $\mathcal{N}_G(h^\star)$ and $\mathcal{N}_G(u)$, the neighbor sets
of $h^\star$ and $u$ respectively. Both are subsets of $V \setminus
\{h^\star, u\}$. $h^\star \notin \mathcal{N}_G(h^\star)$ since there are no
self-loops, and $u \notin \mathcal{N}_G(h^\star)$ by assumption.
Symmetrically, $u \notin \mathcal{N}_G(u)$, and $h^\star \notin
\mathcal{N}_G(u)$ since $u \notin \mathcal{N}_G(h^\star)$ implies $h^\star
\notin \mathcal{N}_G(u)$ on an undirected graph. Hence
\begin{align*}
|\mathcal{N}_G(h^\star)| + |\mathcal{N}_G(u)| &= \deg(h^\star) + \deg(u) \\
&\geq \deg(h^\star) + \delta_{\bar N}(h^\star) \geq N,
\end{align*}
using $\deg(u) \geq \delta_{\bar N}(h^\star)$, since $u$ is by construction
one of the nodes over which $\delta_{\bar N}(h^\star)$ is minimized, and
the theorem's hypothesis.

$|V \setminus \{h^\star, u\}| = N - 2$, and both $\mathcal{N}_G(h^\star)$
and $\mathcal{N}_G(u)$ are subsets of this $(N-2)$-element set. Their sizes
sum to at least $N > N-2$, so they cannot be disjoint: this is a standard
pigeonhole argument, since $|X|+|Y| > |Z|$ for $X,Y\subseteq Z$ forces
$X\cap Y\neq\emptyset$, as otherwise $|X\cup Y| = |X|+|Y| > |Z|$ would
contradict $X\cup Y\subseteq Z$. So there exists $w \in
\mathcal{N}_G(h^\star)\cap\mathcal{N}_G(u)$, giving a length-2 path
$h^\star - w - u$, hence $d(h^\star,u)\leq 2$.
\end{proof}

\begin{theorem}[Depth gap in terms of diameter]
\label{thm:mdst-gap-spectral}
For any connected graph $G$, $D - D' \leq \mathrm{diam}(G)/2$. In the expander
regime, where $\lambda_2$ is bounded away from $0$, $\mathrm{diam}(G) = O(\log N)$,
so $D - D' = O(\log N)$.
\end{theorem}

\begin{proof}
We use two elementary facts about eccentricity that hold for \emph{any}
connected graph and any choice of root.

$D = \mathrm{ecc}_G(h^\star) \leq \mathrm{diam}(G)$. No single vertex's
eccentricity can exceed the graph's overall diameter: the diameter is the
maximum eccentricity over all vertices, so every individual eccentricity,
including $h^\star$'s, is bounded above by it.

$D' = R(G) \geq \mathrm{diam}(G)/2$. Let $c^\star$ achieve the radius, so
$\mathrm{ecc}_G(c^\star) = D' = R(G)$. For any pair $u,v$ realizing the
diameter, $d(u,v)=\mathrm{diam}(G)$, the triangle inequality via $c^\star$
gives $\mathrm{diam}(G) = d(u,v) \leq d(u,c^\star)+d(c^\star,v) \leq
2\,\mathrm{ecc}_G(c^\star) = 2D'$, so $D' \geq \mathrm{diam}(G)/2$.

Combining,
\[
D - D' \leq \mathrm{diam}(G) - \frac{\mathrm{diam}(G)}{2} = \frac{\mathrm{diam}(G)}{2}.
\]
In the expander regime $\lambda_2$ is bounded away from $0$, which forces
$\mathrm{diam}(G) = O(\log N)$ \citep{mohar1991}, giving $D - D' = O(\log N)$.
\end{proof}

\section{Instantiations: Regret and Best Arm Identification}

\begin{corollary}[Regret matches centralized UCB]
\label{cor:regret}
Under the EigenTree relay protocol, the leading term of the expected group
cumulative regret matches a centralized UCB agent on $P$ total pulls, with
$P = \Theta(NT/(2D+1))$.
\end{corollary}

\begin{proof}
Instantiate Theorem~\ref{thm:coordinator} with $b_i$ as agent $i$'s new
per-arm pull counts and reward sums since the last cycle. The hub's
aggregate state $(\hat n^k, \hat s^k)$ for each arm $k$, together with the
total-pull count $P$, is then identical at every decision point to the
state a centralized UCB agent would hold after observing the same pulls.
$P$ itself is maintained exactly: the hub increments it by the number of
uplink contributions received plus its own pull each cycle, which by
Proposition~\ref{prop:aggbcast} is exactly the number of new pulls made
network-wide since the last decision. The standard UCB regret bound is a
function only of $(\hat n^k,\hat s^k,P)$ evaluated at the algorithm's own
internal clock $P$, not wall-clock time, and this sequence of internal
states is identical between the relay protocol and a literal centralized
run. So the leading-order regret term is identical as a function of $P$.

The relationship between $P$ and wall-clock horizon $T$ follows from the
fixed cycle structure. Each cycle lasts $2D+1$ rounds: one simultaneous
pull round, $D$ uplink rounds, one hub-decision round, $D$ downlink
rounds. Every cycle contributes exactly $N$ new pulls, one per agent, to
$P$. Over horizon $T$, the number of completed cycles is
$\Theta(T/(2D+1))$, so $P = \Theta(NT/(2D+1))$.
\end{proof}

\begin{corollary}[BAI sample complexity matches Hillel]
\label{cor:bai}
The total number of arm pulls required by EigenTree-BAI to identify $a^\star$ with probability $\geq 1-\delta$ is identical to Hillel's algorithm; the only overhead is $2D$ extra time steps per elimination round.
\end{corollary}

\begin{proof}
EigenTree-BAI uses exactly Hillel's Algorithm 3 pull schedule $L_r = \lceil t_r
- t_{r-1}\rceil$ and elimination rule, eliminate arms with estimated mean
below the round-$r$ threshold $\epsilon_r$, applied to the hub's aggregate
statistics rather than an all-to-all broadcast. Instantiate
Theorem~\ref{thm:coordinator} with $b_i$ as agent $i$'s per-arm sums over
the current surviving arm set $S_{r-1}$. The hub's aggregate after each
round's uplink is then exactly the sum of pulls Hillel's centralized
algorithm would have observed given the same schedule and elimination
history, since every agent executes the identical, schedule-determined
number of pulls $L_r$ per round. The schedule depends only on
$N,K,\delta,r$, all known in advance, so no agent needs to learn anything
to determine its own pull count. The elimination decision depends only on
this aggregate, so the sequence of surviving sets $S_1 \supseteq S_2
\supseteq \cdots$ produced by EigenTree-BAI is identical to Hillel's, round for
round, and the total number of pulls until $|S_r|=1$ is identical.

The only addition is communication. Each round $r$ requires $D$ uplink
rounds to deliver the round's local sums to the hub, and $D$ downlink
rounds to broadcast the surviving set $S_r$ and the depth signal: $2D$
rounds of overhead per elimination round beyond the $L_r$ pull rounds
themselves.
\end{proof}

\section{Fault Tolerance}
\label{sec:ft}

Throughout this section, $\mathcal{T}$, $D$, $\mathrm{par}(\cdot)$,
$\mathrm{ch}(\cdot)$, $d_i$ refer to the routing tree and its structure
\emph{before} the failure under discussion. A subscript or prime denotes
the corresponding quantity afterward, as introduced at each result.

\begin{definition}[2-connectivity]
\label{def:2conn}
The underlying graph $G$ is \emph{2-connected}: it remains connected after the removal of any single node. Equivalently, $G$ has no cut vertices.
\end{definition}

We take $G$ to be 2-connected throughout, the standing assumption of
Section~\ref{sec:mdst-gap}: it follows from an algebraic connectivity $\lambda_2$
bounded away from $0$ via $\kappa(G)\geq\lambda_2$ \citep{fiedler1973}, and a
degree bound $\delta(G)\geq N/2$ is one simple sufficient condition.

\subsection{Detection and Synchronous Restart}

Before the survivors can re-elect they must detect the failure and enter the new
power iteration in step, with no coordinator and no synchronized clocks. Both
follow from the routing tree already in place. Every node learned its depth $d_i$
and the height $D$ from the first downlink, where depth is the hop count that
message carried, and every ordinary cycle is one uplink-and-downlink traversal of
length $2D+1$ hops. Each node therefore knows, for each of its tree neighbors, the
exact hop of the cycle at which that neighbor's message is due, and runs a single
uniform timer $\tau = 2D+1$ hops against that schedule, suspecting a neighbor dead
when its due message fails to arrive within $\tau$. The rule is identical for
every node, and no node is a distinguished detector: the hub is detected by its
children exactly as any node is detected by its parent.

\begin{proposition}[Failure horizon]
\label{prop:horizon}
Under the uniform timer $\tau = 2D+1$, after any single node failure every surviving node receives the failure notice within $2D$ hops of the earliest detection.
\end{proposition}

\begin{proof}
Removing $v$ severs $\mathcal{T}$ into components: the component above $v$ (still
containing every node not below $v$) and one component per child
$c\in\mathrm{ch}(v)$, since every tree path between two of these components ran
through $v$. No survivor can flood all of them over tree edges, because the paths
between components are exactly the ones that passed through $v$; instead each
component has its own first detector, the parent $\mathrm{par}(v)$ for the upper
component and the child $c$ for each lower one, and each detector floods the
notice over the intact tree edges of its own component only.

Two quantities bound the horizon. First, the detectors do not fire
simultaneously: a detector at depth $d$ expects the failed node's message at hop
$d$ (uplink) or hop $2D-d$ (downlink) of the cycle, so two detectors' firing
times differ by their difference in these expected hops, which is at most the
depth spread $D$. Since all detectors use the same $\tau$ measured against the
same round-1 schedule, the last detector fires at most $D$ hops after the first.
Second, within a component the notice floods over tree edges from the detector,
reaching the farthest node in at most the component's height, itself at most $D$
hops. Adding the two, the last node to hear does so at most $D + D = 2D$ hops
after the earliest detection.
\end{proof}

The same $2D$ budget makes the restart synchronous. The notice is flooded
carrying the downlink's hop counter, so a node that receives it $h$ hops from its
detector knows that $h$ of the $2D$ hops have elapsed and waits the remaining
$2D-h$ hops before taking its first power-iteration step. Every node's
$(\text{hops travelled}) + (\text{hops waited}) = 2D$, so all survivors take
their first step at the same instant, measured from the earliest detection, with
no clock agreement. The wait is read directly off the hop counter, which is the
node's depth $d_i$ up to the fixed offset it always was, so the depth each node
already holds serves three roles at once, the detection deadline through $\tau$,
the bound on detector skew, and the restart countdown, and no field is added to
any message.

\subsection{Re-Election}

Deleting node $v$ from $A$ gives $\tilde A = A + E$ with
$E = -(A e_v e_v^\top + e_v e_v^\top A)$. As $A$ has no self-loops, $E$ has rank
at most $2$ and $\|E\|_2 = \|A e_v\|_2 = \sqrt{\deg(v)}$, since $A e_v$ is the
$v$-th column of $A$, a $0/1$ vector with $\deg(v)$ ones. This holds for any $v$;
the hub is only the case of largest $\deg(v)$.

\begin{proposition}[Warm-started re-election]
\label{prop:reelection}
After any single node failure, the warm-started power iteration on $\tilde A$ converges to the leading eigenvector of $\tilde A$ from the pre-failure $\tilde\psi$. After the fixed budget $\tau_{\mathrm{init}} = O(\log N)$, the residual angle is at most $(\lambda_2/\lambda_1)^{\tau_{\mathrm{init}}}\cdot\sqrt{\deg(v)}/\gamma$, where $\gamma$ is the eigenvalue gap of $\tilde A$, against $(\lambda_2/\lambda_1)^{\tau_{\mathrm{init}}}\cdot\Theta(1)$ for a cold start from uniform. The warm start is therefore never worse, and its advantage is a factor $\gamma/\sqrt{\deg(v)}$, largest when the failed node has low degree.
\end{proposition}

\begin{proof}
After $v$ goes silent, a surviving node's update
$\tilde w_i = \sum_{j\sim i} w_j$ no longer includes the dead neighbor, so the
operator applied each round is exactly $\tilde A$ on the live nodes. Since $G$ is
2-connected, $G\setminus\{v\}$ is connected, so $\tilde A$ is irreducible and has
a unique positive leading eigenvector $\tilde\psi_{\tilde A}$ by
Perron--Frobenius, and power iteration converges to it from any start with
nonzero overlap; the warm start $\tilde\psi$ qualifies, as both vectors are
strictly positive.

For the accuracy, the initial angle to the target is bounded by Davis--Kahan
\citep{daviskahan1970},
$\sin\Theta(\tilde\psi,\tilde\psi_{\tilde A}) \leq \|E\|_2/\gamma
= \sqrt{\deg(v)}/\gamma$, using $\|E\|_2 = \sqrt{\deg(v)}$ for the rank-2
perturbation $E$ deleting row and column $v$. Each power-iteration round
multiplies the component orthogonal to $\tilde\psi_{\tilde A}$ by at most
$\lambda_2/\lambda_1$ relative to the leading one, so the residual angle after
$t$ rounds is at most $(\lambda_2/\lambda_1)^t\tan\Theta$.

The survivors run the fixed budget $t = \tau_{\mathrm{init}} = O(\log N)$
rather than a tolerance-driven number of rounds. This is forced: a budget
derived from the bound above would require $\deg(v)$, which only $v$'s
neighbors know and which $v$ cannot report once dead, and $\gamma$, a spectral
quantity of $\tilde A$ that no node holds. Every node can compute
$\tau_{\mathrm{init}}$ from $N$ alone, so all survivors finish on the same
round, exactly as at initialization.

Substituting $t = \tau_{\mathrm{init}}$ and the two initial angles gives the
claim: the warm start begins at $\tan\Theta \leq \sqrt{\deg(v)}/\gamma$, while
a cold start from the uniform vector begins at $\tan\Theta = \Theta(1)$, having
no dependence on the failure. Both contract at the same rate over the same
budget, so the warm start's residual is smaller by the factor
$\gamma/\sqrt{\deg(v)}$, and it is never worse. The advantage is largest when
$\deg(v)$ is small; when $\deg(v) = \Theta(N)$ it narrows to
$\gamma/\sqrt{N}$, since $\deg(v) \leq N$ always.
\end{proof}

\begin{remark}
\label{rem:gamma}
The bound treats $\gamma$, the eigenvalue gap of the \emph{perturbed} matrix
$\tilde A$, as $\Omega(1)$ in the expander regime. This is a property of
$\tilde A$, not of $A$, so it requires the gap to survive the failure. Since
$E$ has rank at most $2$, Weyl's inequality moves each eigenvalue of $A$ by at
most $\|E\|_2 = \sqrt{\deg(v)}$, which bounds the perturbed gap in terms of the
original one but does not by itself keep it away from $0$ when $\deg(v)$ is
large. A failure-independent lower bound on $\gamma$ is left for future work.
\end{remark}

\subsection{Recovery}

A recovering node $v$, previously failed and now functioning again, has no
reliable state: it may have missed any number of promotions and
re-routings during its absence. Rather than attempting to recover its old
position in $\mathcal{T}$, it rejoins as if it were new. It sends a
\textsc{Ping} to each of its graph-neighbors $\mathcal{N}_G(v)$, and any
live neighbor $j$ that is itself already attached to the current tree
replies with its current depth $d_j$. $v$ attaches as a child of whichever
neighbor replies with the smallest $d_j$, ties broken by smallest node
identity, mirroring the tie-break already used during initial
max-flooding. It sets $d_v \leftarrow d_j+1$ and waits for the next
downlink through $j$ to populate its mirrored aggregate state, exactly as
any other node does at a cycle boundary.

\begin{proposition}[Recovery correctness]
\label{prop:recovery-correct}
If $v$ attaches to a live, already-attached neighbor $j$ upon recovery, the resulting structure remains a valid spanning tree over all currently live nodes, with no cycles introduced, regardless of which node currently serves as hub.
\end{proposition}

\begin{proof}
By hypothesis, immediately before $v$'s recovery the live nodes other than
$v$ form a valid spanning tree $\mathcal{T}''$ rooted at whichever node is
currently hub: the original $h^\star$ or a hub elected after a later failure.
The argument below does not depend on which. $v$ itself is disconnected from
$\mathcal{T}''$, having no parent pointer and no place in the current
tree.

The recovery procedure adds exactly one new edge, $v \to j$, where $j$ is
a node already in $\mathcal{T}''$: by construction, $v$ only attaches to a
neighbor that replies to its ping, and only already-attached live
neighbors reply. $v$ was not previously part of $\mathcal{T}''$ at all, so
this single new edge cannot close a cycle. Any cycle would require a
second path between $v$ and some node of $\mathcal{T}''$, but $v$ had no
edges into $\mathcal{T}''$ before this step. The resulting graph is
therefore $\mathcal{T}''$ together with one pendant edge attaching $v$ as
a new leaf child of $j$, which is again a tree, since a tree plus a
pendant edge to a new vertex is a tree, spanning all live nodes including
$v$, rooted at the same node as $\mathcal{T}''$.

This argument uses only that $j$ is live and already attached. It does not
depend on $j$'s identity, depth, or role, whether ordinary node or current
hub, so it composes without modification with any prior sequence of
re-elections.
\end{proof}

\bibliography{references}

\end{document}
